%@ pages 242-250
%@ status translated
% Bibliography (original pp. 242--250).
% Bibliographic titles are left in the language in which they were printed;
% only the authors' annotations are translated.  Two obvious slips are
% corrected silently in the reference list (Y. L. "Ton." -> "Tong";
% "Grundelehren" -> "Grundlehren").

\begin{thebibliography}{Hopf}
\phantomsection\label{references}

\bibitem[xxx]{xxx}
For the alg-geom, q-alg and math.AT preprints see
\newline
{\tt http://xxx.lanl.gov/form/math}
\newline (you should first look for the nearest mirror site, see \newline
{\tt http://xxx.lanl.gov/servers.html}).

\bibitem[Hopf]{Hopf}
For the Hopf preprints see \newline
{\tt http://hopf.math.purdue.edu/pub/new-html/cgi-interface.html}
\newline
(search by author).

\bibitem{Adams}
J. Adams. {\em Infinite Loop Spaces}, Princeton University Press {\em Annals
of Math. Studies} {\bf 90} (1978).

\bibitem{AlbertMolino}
C. Albert, P. Molino. {\em Pseudogroupes de Lie transitifs, II:
th\'eor\`emes d'int\'egrabilit\'e}, Hermann, Paris (1987).

\bibitem{SGA4}
M. Artin, J. Verdier, A. Grothendieck. {\em SGA 4: Th\'eorie des Topos et
Cohomologie Etale des Sch\'emas} Springer Lecture Notes in Mathematics {\bf
269},  Heidelberg (1972-1973).

\bibitem{SGA4b}
M. Artin, J. Verdier, A. Grothendieck. {\em SGA 4: Th\'eorie des Topos et
Cohomologie Etale des Sch\'emas} Springer Lecture Notes in Mathematics
{\bf 270},  Heidelberg (1972-1973).

\bibitem{SGA4c}
M. Artin, J. Verdier, A. Grothendieck. {\em SGA 4: Th\'eorie des Topos et
Cohomologie Etale des Sch\'emas} Springer Lecture Notes in Mathematics {\bf
305}, Heidelberg
(1972-1973).

\bibitem{ArtinInventiones}
M. Artin.  Versal deformations and algebraic stacks, {\em Inventiones Math.}
{\bf 27} (1974), 165-189.

\bibitem{BaezDolanLetter}
J. Baez, J. Dolan. $n$-Categories, sketch of a
definition. Letter to R. Street, 29 Nov. and 3 Dec. 1995, available at
{\tt http://math.ucr.edu/home/baez/ncat.def.html}

\bibitem{BaezDolan}
J. Baez, J. Dolan. Higher-dimensional algebra and
topological quantum field theory. {\em  Jour. Math. Phys} {\bf  36} (1995),
6073-6105.

\bibitem{BaezDolanIII}
J. Baez, J. Dolan. Higher dimensional algebra III: $n$-categories and the
algebra of opetopes.
Preprint
q-alg 9702014, to appear in {\em Adv. Math.}.

\bibitem{Batanin}
M. Batanin. On the definition of weak $\omega$-category. Macquarie mathematics
report number 96/207, Macquarie University, NSW Australia.

\bibitem{Benabou}
J. B\'enabou. {\em Introduction to Bicategories}, Lect. Notes in Math. {\bf
47}, Springer-Verlag (1967).

\bibitem{SGA6}
P. Berthelot, A. Grothendieck, L. Illusie. {\em S.G.A. 6:
Th\'eorie des Intersections et Th\'eor\`eme de Riemann-Roch} Springer Lecture
Notes in Mathematics {\bf 225}, Heidelberg (1971).

\bibitem{BoardmanVogt}
J. Boardman, R. Vogt. {\em Homotopy invariant algebraic structures
on topological spaces} Springer Lecture Notes in Mathematics {\bf 347}, Berlin
(1973).

\bibitem{BondalKapranov}
A. J. Bondal, M.M. Kapranov. Enhanced triangulated categories.
{\em Math. USSR Sb.},
{\bf 70} (1991) 93-107.

\bibitem{BousfieldKan}
A. Bousfield, D. Kan. {\em Homotopy limits, completions and localisations}.
Lecture Notes in Math. {\bf 304}, Springer-Verlag (1972).

\bibitem{BousfieldFriedlander}
A. Bousfield, E. Friedlander. Homotopy theory of $\Gamma$-spaces, spectra and
bisimplicial sets. {\em Lecture Notes in Math.} {\bf 658}, Springer (1978),
80-150.

\bibitem{BreenIHES}
L. Breen. Extensions du groupe additif. {\em Publ. Math. I.H.E.S.} {\bf 48}
(1978), 39-126.

\bibitem{Breen}
L. Breen. On the classification of $2$-gerbs and $2$-stacks. {\em Ast\'erisque}
{\bf 225}, Soc. Math. de France (1994).

\bibitem{KBrown}
K. Brown. Abstract homotopy theory and generalized sheaf cohomology. {\em
Trans. A.M.S.} {\bf 186} (1973), 419-458.

\bibitem{BrownGersten}
K. Brown, S. Gersten. {\em Algebraic $K$-theory as generalized sheaf cohomology}
Springer {\em Lecture Notes in Math.} {\bf 341} (1973), 266-292.

\bibitem{RBrown}
R. Brown. The twisted Eilenberg-Zilber theorem, {\em Celebrazioni Archimedee del
secolo XX, Simposio di topologia} (1967), 34-37.

\bibitem{Chai-Faltings}
C.-L. Chai, G. Faltings. {\em Degenerations of Abelian Varieties},
Springer-Verlag (1990).

\bibitem{Cobb}
P. Cobb. $P_n$-spaces and $n$-fold loop spaces. {\em Bull. Amer. Math. Soc.}
{\bf 80} (1974), 910-914.

\bibitem{CordierPorter}
J.-M. Cordier, T. Porter. Homotopy-coherent category theory. {\em Trans. Amer.
Math. Soc.} {\bf 349} (1997), 1-54.

\bibitem{Deligne-Mumford}
P. Deligne, D. Mumford. The irreducibility of the space of curves of a given
genus, {\em  Publ. Math. I.H.E.S.} {\bf 36}, (1969), 75-110.

\bibitem{HodgeIII}
P. Deligne. Th\'eorie de Hodge III. {\em  Publ. Math. I.H.E.S.}
{\bf 44} (1974), 5-78.

\bibitem{DoldPuppe}
A. Dold, D. Puppe. Homologie nicht-additiver Funktoren. Anwendungen. {\em
Ann. Inst. Fourier} {\bf
11} (1961), 201-312.

\bibitem{Dunn}
G. Dunn. Uniqueness of $n$-fold delooping machines. {\em J. Pure and Appl. Alg.}
{\bf 113} (1996), 159-193.

\bibitem{Dunn97}
G. Dunn. Lax operad actions and coherence for monoidal $n$-categories,
$A_{\infty}$ rings and modules. {\em Theory and applications of categories}
{\bf 3} (1997), 50-84.

\bibitem{DHK}
W. Dwyer, P. Hirschhorn, D. Kan. Model categories and more general abstract
homotopy theory: a work in what we like to think of as progress. Preprint,
available at {\tt http://www-math.mit.edu/~psh/}.

\bibitem{DwyerKan1}
W. Dwyer, D. Kan. Simplicial localizations of categories. {\em J. Pure and Appl.
Algebra} {\bf 17} (1980), 267-284.

\bibitem{DwyerKan2}
W. Dwyer, D. Kan. Calculating simplicial localizations. {\em J. Pure and
Appl. Algebra} {\bf 18} (1980), 17-35.

\bibitem{DwyerKan3}
W. Dwyer, D. Kan. Function complexes in homotopical algebra. {\em Topology}
{\bf 19} (1980), 427-440.

\bibitem{DwyerKanDiags}
W. Dwyer, D. Kan. Equivalences between homotopy theories of diagrams.
{\em Algebraic Topology and Algebraic $K$-theory}, {\em Annals of Math. Studies}
{\bf 113}, Princeton University Press (1987), 180-205.

\bibitem{DwyerE2}
W. Dwyer, D. Kan, C. Stover. An $E^2$ model category structure for pointed
simplicial spaces. {\em J. Pure and Appl. Alg.} {\bf 90} (1993), 137-152.

\bibitem{Ehresmann}
C. Ehresmann. Structures feuillet\'ees. {\em Proc. of the 5th
Canadian Math. Cong.}

\bibitem{Ehresmann2}
C. Ehresmann. Structures locales et structures infinit\'esimales.
{\em C.R.A.S.} {\bf 234} (1952), 587-589.

\bibitem{Ehresmann3}
C. Ehresmann. Cat\'egories topologiques et cat\'egories diff\'erentiables.
{\em Coll. G\'eom. Diff. Globale Bruxelles} (1958), 137-150.

\bibitem{Ehresmann4}
C. Ehresmann. Introduction \`a la th\'eorie des structures
infinit\'esimales et des
pseudogroupes de Lie. {\em Coll. Int. CNRS} Strasbourg (1959).

\bibitem{Fiedorowicz}
Z. Fiedorowicz. Classifying spaces of topological monoids and categories.
{\em Amer. J. Math.} {\bf 106} (1984), 301-350.

\bibitem{GabrielZisman}
P. Gabriel, M. Zisman. {\em Calculus of fractions and homotopy theory},
Ergebnisse der Math. und ihrer Grenzgebiete {\bf 35}, Springer-Verlag, New York
(1967).

\bibitem{Giraud}
J. Giraud. {\em Cohomologie nonab\'elienne}, Die Grundlehren der
Math. Wissen\-schaften {\bf 179} Springer-Verlag (1971).

\bibitem{GiraudThese}
J. Giraud. {\em M\'ethode de la descente.} Memoire {\bf 2} de la S.M.F.
(1964).

\bibitem{Godbillon}
C. Godbillon. {\em Feuilletages: \'etudes g\'eom\'etriques}
Birkh\"auser Progress in Math. {\bf 98}, Boston (1991).

\bibitem{Godement}
R. Godement. {\em Topologie alg\'ebrique et th\'eorie des faisceaux.}
Hermann, Paris (1958).

\bibitem{JardineGoerssBook}
P. Goerss, J. Jardine. {\em Simplicial homotopy theory}, book available
as a preprint at {\tt
http://www.math.uwo.ca/}$\sim${\tt jardine/papers/simp-sets.}

\bibitem{JardineGoerssLoc}
P. Goerss, J. Jardine. Localization theories for simplicial presheaves. Preprint
available at the Hopf archive.

\bibitem{Gordon-Power-Street}
R. Gordon, A.J. Power, R. Street.  Coherence for tricategories {\em  Memoirs
A.M.S.} {\bf 117} (1995), 558 ff.

\bibitem{Grothendieck}
A. Grothendieck.  {\em Pursuing Stacks}, unpublished, available from the
Universit\'e de Montpellier 2 (J. Malgoire), or the University of Wales
(R. Brown).

\bibitem{GrothendieckDix}
A. Grothendieck.  {\em Dix expos\'es sur la cohomologie des sch\'emas},
North-Holland, Amsterdam (1968).

\bibitem{SGA1}
A. Grothendieck. {\em Rev\^etements \'etales et groupe fondamental}
Springer Lecture Notes in Mathematics {\bf 224} Heidelberg (1971).

\bibitem{Gugenheim}
V.K.A.M. Gugenheim. On a chain complex of a fibration. {\em Illinois J. Math.}
{\bf 3} (1972), 398-414.

\bibitem{GugenheimLambe}
V.K.A.M. Gugenheim, L. Lambe. Applications of perturbation theory to
differential homological algebra I,II. {\em Illinois J. Math.} {\bf 33} (1989),
556-582; {\bf 35} (1991), 357-373.

\bibitem{Haefliger}
A. Haefliger. Structures feuillet\'ees et cohomologie \`a valeurs dans
un faisceau de groupo\"{\i}des. {\em Comment. Math. Helv.} {\bf 32} (1958),
215-223.

\bibitem{Haefliger2}
A. Haefliger. Groupo\"{\i}des de holonomie et classifiants. {\em
Ast\'erisque} {\bf 116} (1984),
70-97.

\bibitem{ResiduesDuality}
R. Hartshorne. {\em Residues and Duality} Springer {\em L.N.M.} {\bf 20}
(1966).

\bibitem{Hinich1}
V. Hinich. Descent of Deligne groupoids. Preprint alg-geom 9606010.

\bibitem{Hinich2}
V. Hinich. Homological algebra of homotopy algebras. Preprint q-alg 9702015.

\bibitem{Hirschhorn}
P. Hirschhorn. {\em Localization of model categories}
book in preparation; available at
{\tt http://www-math.mit.edu/~psh/}.

\bibitem{HoveyBook}
M. Hovey. {\em Model categories.} Book preprint, 1997, available
at Hopf.

\bibitem{HoveyMonMod}
M. Hovey. Monoidal model categories. Preprint, 1998, available at the
Hopf archive,
also math.AT 9803002.

\bibitem{Illusie}
L. Illusie. {\em Le complexe cotangent I et II}, {\em  Lecture Notes in
Math.} {\bf 239, 283}, Springer-Verlag, Berlin (1971,1972).

\bibitem{Jardine}
J.F. Jardine.  Simplicial presheaves, {\em J. Pure and Appl. Algebra} {\bf 47}
(1987), 35-87.

\bibitem{JardineBool}
J. Jardine. Boolean localization, in practice. {\em Doc. Math.} {\bf 1} (1996),
245-275. This journal is available at
{\tt http://www.mathematik.uni-bielefeld.de/documenta/}.

\bibitem{Joyal}
A. Joyal, Letter to A. Grothendieck (cf.\ \cite{Jardine}).

\bibitem{Kelly}
G. Kelly, {\em Basic concepts of enriched category theory} London Math. Soc.
Lecture Notes {\bf 64}, Cambridge U. Press, Cambridge (1982).

\bibitem{Kontsevich}
M. M. Kontsevich. Homological algebra of mirror symmetry. Proc. ICM 94.
Birkh\"auser-Verlag 1995, 120-139.

\bibitem{LaumonMB}
G. Laumon, L. Moret-Bailly. Champs alg\'ebriques. Preprint, Orsay {\bf 42}
(1992).

\bibitem{Leroy}
O. Leroy.  Sur une notion de $3$-cat\'egorie adapt\'ee \`a l'homotopie.
Preprint, Univ. de Montpellier 2 (1994).

\bibitem{MacLane}
S. Mac Lane. {\em Categories for the working mathematician.}
Springer, Berlin (1988).

\bibitem{MacLane2}
S. Mac Lane. Natural associativity and commutativity. {\em Rice U. Studies}
{\bf 49} (1963), 28-46.

\bibitem{MacLaneMoerdijk}
S. Mac Lane, I. Moerdijk. {\em Sheaves in Geometry and Logic: A First
Introduction to Topos Theory.} Springer-Verlag, Heidelberg (1992).

\bibitem{Mochizuki}
S. Mochizuki.  The geometry of the compactification  of the Hurwitz scheme.
{\em Publ. R.I.M.S. Kyoto University} {\bf 31} (1995), 355-441.

\bibitem{Moerdijk}
I. Moerdijk. On the weak homotopy type of etale groupoids. {\em Progress
in Math.} {\bf 145}
({\em Integrable systems and foliations}, C. Albert {\em et al} eds.)
Birkh\"auser, Boston (1997),
147-156.

\bibitem{Molino}
P. Molino. {\em Riemannian foliations.} Birkh\"auser, Boston (1988).

\bibitem{Morita}
S. Morita. On characteristic classes of Riemannian foliations. {\em
Osaka Math. Jour.} {\bf 16} (1979), 161-172.

\bibitem{ObrianToledoTong1}
N. R. O'Brian, D. Toledo, Y. L. Tong.  Hirzebruch-Riemann-Roch for
coherent sheaves, {\em Amer. Jour. Math.} {\bf  103} (1981),  253-271.

\bibitem{ObrianToledoTong2}
N. R. O'Brian, D. Toledo,  Y. L. Tong. Grothendieck-Riemann-Roch for
analytic maps.  {\em Math Annalen} {\bf  271} (1985), 493-526.

\bibitem{Pradines}
J. Pradines. Th\'eorie de Lie pour les groupo\"{\i}des
diff\'erentiables. Relations entre propri\'et\'es locales et globales.
{\em C.R.A.S.} {\bf 263} (1966), 907-910.

\bibitem{Pradines2}
J. Pradines. Th\'eorie de Lie pour les groupo\"{\i}des
diff\'erentiables. Calcul diff\'erentiel dans la cat\'egorie des
vari\'et\'es diff\'erentiables. {\em C.R.A.S.} {\bf 264} (1967),
245-248.

\bibitem{Pradines3}
J. Pradines. G\'eom\'etrie diff\'erentielle au-dessus d'un
groupo\"{\i}de. {\em C.R.A.S.} {\bf 266} (1968), 1194-1196.

\bibitem{Pradines4}
J. Pradines. Troisi\`eme th\'eor\`eme de Lie pour les groupo\"{\i}des
diff\'erentiables.
{\em C.R.A.S.} {\bf 267} (1968), 21-23.

\bibitem{Quillen}
D. Quillen. {\em Homotopical algebra.} Springer, {\em L.N.M.} {\bf 43} (1967).

\bibitem{QuillenAnnals}
D. Quillen. Rational homotopy theory. {\em Ann. of Math.} {\bf 90} (1969),
205-295.

\bibitem{QuillenK}
D. Quillen. Higher algebraic K-theory I, {\em Higher K-theories} Springer {\em
L.N.M.} {\bf 341} (1973), 85-147.

\bibitem{Reedy}
C. L. Reedy. Homotopy theory of model categories. Unpublished preprint (1974),
recently reissued by P. Hirschhorn and available at \newline
{\tt http://www-math.mit.edu/~psh/}

\bibitem{Rezk}
C. Rezk. Email to P. Hirschhorn, 4 July 1998.

\bibitem{Segal}
G. Segal. Homotopy everything $H$-spaces. Preprint.

\bibitem{SegalInventiones}
G. Segal. Configuration spaces and iterated loop spaces. {\em Inv. Math.} {\bf
21} (1973), 213-221.

\bibitem{SegalTopology}
G. Segal. Categories and cohomology theories. {\em Topology} {\bf 13}
(1974), 293-312.

\bibitem{flexible}
C. Simpson. Flexible sheaves. Preprint, q-alg 9608025.

\bibitem{kobe}
C. Simpson. Homotopy over the complex numbers and generalized de Rham
cohomology. {\em Moduli of Vector Bundles}, M. Maruyama (Ed.) {\em Lecture
Notes in Pure and Applied Math.} {\bf 179}, Marcel Dekker (1996), 229-263.

\bibitem{geometricN}
C. Simpson. Algebraic (geometric) $n$-stacks. Preprint alg-geom
96-09014.

\bibitem{nCAT}
C. Simpson. A closed model structure for $n$-categories, internal $Hom$,
$n$-stacks and generalized Seifert-Van Kampen. Preprint alg-geom 9704006.

\bibitem{limits}
C. Simpson. Limits in $n$-categories. Preprint alg-geom 9708010.

\bibitem{effective}
C. Simpson. Effective generalized Seifert-Van Kampen.
Preprint q-alg 9710011.

\bibitem{Stasheff}
J. Stasheff. Homotopy associativity of $H$-spaces. {\em Trans. Amer. Math.
Soc.} {\bf 108} (1963), 275-312.

\bibitem{Tamsamani}
Z. Tamsamani.  Sur des notions de $n$-cat\'egorie et $n$-groupo\"{\i}de non-stricte
via des ensembles multi-simpliciaux. Thesis, Universit\'e Paul Sabatier,
Toulouse (1996) available on alg-geom (95-12 and 96-07).

\bibitem{ThomasonLimits}
R. Thomason. Homotopy colimits in the category of small categories. {\em Math.
Proc. Cam. Phil. Soc.} {\bf 85} (1979), 91-110.

\bibitem{Thomason}
R. Thomason.  Algebraic K-theory and etale cohomology, {\em Ann. Sci. E.N.S.}
{\bf 18} (1985), 437-552.

\bibitem{ThomasonCatMod}
R. Thomason. {\bf Cat} as a closed model category. {\em Cahiers Top.
G\'eom. Diff.} {\bf 21}
(1980), 305-324.

\bibitem{ThomasonTrobaugh}
R. Thomason, T. Trobaugh. Higher algebraic $K$-theory of schemes and of
derived categories.
{\em The Grothendieck Festschrift, Vol. III} Birkh\"auser, Boston (1990),
247-436.

\bibitem{ToenDescente}
B. Toen. Un th\'eor\`eme de descente. Note, Universit\'e de
Toulouse 3 (1996).

\bibitem{Toledo1}
D. Toledo,  Y. L. Tong., A parametrix for {d-bar}  and Riemann-Roch in Cech
theory, {\em Topology}, {\bf 15} (1976),  273-301.

\bibitem{Toledo2}
D. Toledo, Y. L. Tong. Duality and intersection theory in complex
manifolds, II:  The holomorphic Lefschetz formula, {\em Annals of Math.} {\bf
108} (1978).

\bibitem{Toledo3}
D. Toledo, Y. L. Tong. Green's theory of Chern classes and the
Riemann-Roch formula, {\em  Contemporary Mathematics} Vol. 58, Part I, 261-275.

\bibitem{TravauxThomason}
C. Wiebel. The mathematical enterprises of Robert Thomason. {\em
Bulletin A. M. S.} {\bf 34} (1997) 1-13.

Here are some references added in version 2:

\bibitem{Cordier}
J. M. Cordier. Comparaison de deux cat\'egories d'homotopie de morphismes
coh\'erents. {\em Cahiers Top. G\'eom. Diff. Cat.} {\bf 30} (1989),
257-275.

\bibitem{CordierPorter2}
J. M. Cordier, T. Porter. Fibrant diagrams, rectifications and a construction
of Loday. {\em J. P. A. A.} {\bf 67} (1990), 111-124.

\bibitem{CordierPorter3}
J. M. Cordier, T. Porter. Vogt's theorem on categories of homotopy-coherent
diagrams. {\em Math. Proc. Camb. Phil. Soc.} {\bf 100} (1986), 65-90.

\bibitem{DKS}
W. Dwyer, D. Kan, J. Smith. Homotopy-commutative diagrams and their
realizations. {\em J. P. A. A.} {\bf 57} (1989), 5-24.

\bibitem{EH}
D. Edwards, H. Hastings. \v{C}ech and Steenrod homotopy theories with
applications to geometric topology. {\em L. N. M.} {\bf 542},
Springer (1976).

\bibitem{Leitch}
Leitch. The homotopy-commutative cube. {\em J. London Math. Soc.}
{\bf 9} (1974/75), 23-29.

\bibitem{Mather}
M. Mather. Pull-backs in homotopy theory. {\em Canadian J. Math.} {\bf 28}
(1976), 225-263.

\bibitem{Vogt}
R. Vogt. Homotopy limits and colimits. {\em Math. Z.} {\bf 134}
(1973), 11-52.

\bibitem{Vogt2}
R. Vogt. Commuting homotopy limits. {\em Math. Z.} {\bf 153} (1977), 59-82.

Reference for {\tt v3}:

\bibitem{Dugger}
D. Dugger. Simplicial presheaves, revisited. In preparation.

\end{thebibliography}
